\documentclass[../../00_Main/Main.tex]{subfiles}
\begin{document}
The main goal of this work is the design and development of a \gls{phantom} 
for the evaluation and calibration of \acrfull{eeg} signal acquisition systems.  
To achieve this goal, the following specific objectives are proposed:

\begin{enumerate}

\item Design a \gls{phantom} capable of simulating realistic conditions 
for the transmission of \acrshort{eeg} signals \nota{(\textbf{and other standard test signals, like sinusoids and square waves})} through a physical medium that mimics, 
through an internal generation system that maintains stable operation for at least 
\nota{X (to be defined) days/hours - 1 day ?}.

\item Develop a desktop application for operational integration with the \nota{designed/proposed} hardware system, 
intended for the comprehensive management of simulated \acrshort{eeg} signals within the \gls{phantom}.  
The platform should serve as a unified working environment, offering key functionalities such as:

\begin{enumerate}
    \item Control of signal generation parameters, with flexible configuration of amplitude, frequency, and waveform.  
    \nota{this “waveform” aspect could mean either uploading a file with a predefined structure (ex: csv with some predefined column names) to reproduce that waveform, 
    and/or allowing the user to select within the application options such as “generate sinusoidal wave with frequency $f$ and amplitude $A$.” 
    We will later refine what is actually permitted and phrase this section accordingly.}
    \nota{Approx.\ –50 to 50~mV are the expected target values, covering both EEG amplitudes and some useful synthetic test cases with sinusoids or other test functions.}

    \item Real-time visualization of the acquired signals (with minimal latency and jitter), 
    including graphical tools for monitoring and validation.  
    \nota{we should define a numerical cutoff for acceptable latency and jitter — e.g., \textbf{up to} 20~ms latency 
    and 5~ms jitter — since it reads more clearly as an objective than “as low as possible.”}

    \item Structured data logging, with the possibility of exporting recorded measurements 
    to other environments for further analysis and processing.  
    \nota{we could add here something like: data recording and export in format~X, 
    which is widely accepted by other analysis platforms.}
\end{enumerate}

\item Ensure system compatibility with the \textit{OpenBCI} platform by adopting communication standards 
and data structures that allow correct reception and processing of generated signals.  
This integration with open hardware aims to guarantee system interoperability, facilitate replicability 
in other research contexts, and promote adoption, adaptation, and improvement by third parties 
within an open development framework.  
\nota{Chequear redacción de este item.}

\end{enumerate}


\biblio % Needed for referencing to working when compiling individual subfiles - Do not remove
\end{document}
